%%% The main file. It contains definitions of basic parameters and includes all other parts.

% Meta-data of your thesis (please edit)
\input metadata.tex

% Generate metadata in XMP format for use by the pdfx package
\input xmp.tex

%% Settings for single-side (simplex) printing
% Margins: left 40mm, right 25mm, top and bottom 25mm
% (but beware, LaTeX adds 1in implicitly)
\documentclass[12pt,a4paper]{report}
\setlength\textwidth{145mm}
\setlength\textheight{247mm}
\setlength\oddsidemargin{15mm}
\setlength\evensidemargin{15mm}
\setlength\topmargin{0mm}
\setlength\headsep{0mm}
\setlength\headheight{0mm}
% \openright makes the following text appear on a right-hand page
\let\openright=\clearpage

%% Settings for two-sided (duplex) printing
% \documentclass[12pt,a4paper,twoside,openright]{report}
% \setlength\textwidth{145mm}
% \setlength\textheight{247mm}
% \setlength\oddsidemargin{14.2mm}
% \setlength\evensidemargin{0mm}
% \setlength\topmargin{0mm}
% \setlength\headsep{0mm}
% \setlength\headheight{0mm}
% \let\openright=\cleardoublepage

%% If the thesis has no printed version, symmetric margins look better
% \documentclass[12pt,a4paper]{report}
% \setlength\textwidth{145mm}
% \setlength\textheight{247mm}
% \setlength\oddsidemargin{10mm}
% \setlength\evensidemargin{10mm}
% \setlength\topmargin{0mm}
% \setlength\headsep{0mm}
% \setlength\headheight{0mm}
% \let\openright=\clearpage

%% Generate PDF/A-2u
\usepackage[a-2u]{pdfx}

%% Prefer Latin Modern fonts
\usepackage{lmodern}
% If we are not using LuaTeX, we need to set up character encoding:
\usepackage{iftex}
\ifpdftex
\usepackage[utf8]{inputenc}
\usepackage[T1]{fontenc}
\usepackage{textcomp}
\fi

%% Further useful packages (included in most LaTeX distributions)
\usepackage{amsmath}        % extensions for typesetting of math
\usepackage{amsfonts}       % math fonts
\usepackage{amsthm}         % theorems, definitions, etc.
\usepackage{bm}             % boldface symbols (\bm)
\usepackage{booktabs}       % improved horizontal lines in tables
\usepackage{caption}        % custom captions of floating objects
\usepackage{dcolumn}        % improved alignment of table columns
\usepackage{floatrow}       % custom float environments
\usepackage{graphicx}       % embedding of pictures
\usepackage{indentfirst}    % indent the first paragraph of a chapter
\usepackage[nopatch=item]{microtype}   % micro-typographic refinement
\usepackage{paralist}       % improved enumerate and itemize
\usepackage[nottoc]{tocbibind} % makes sure that bibliography and the lists
			    % of figures/tables are included in the table
			    % of contents
\usepackage{xcolor}         % typesetting in color

% The hyperref package for clickable links in PDF and also for storing
% metadata to PDF (including the table of contents).
% Most settings are pre-set by the pdfx package.
\hypersetup{unicode}
\hypersetup{breaklinks=true}

% Packages for computer science theses
\usepackage{algpseudocode}  % part of algorithmicx package
\usepackage{algorithm}
\usepackage{fancyvrb}       % improved verbatim environment
\usepackage{listings}       % pretty-printer of source code

% Custom packages
\usepackage{enumitem}
\usepackage{tikz}
\usepackage{forest}
\usepackage{siunitx}
\usepackage{longtable}
\usepackage{pdflscape}
\usepackage{adjustbox}

\usepackage{booktabs, multirow} % for borders and merged ranges
\usepackage{soul}% for underlines
\usepackage{xcolor,colortbl} % for cell colors
\usepackage{changepage,threeparttable} % for wide tables

% You might want to use cleveref for references
% \usepackage{cleveref}

% Set up formatting of bibliography (references to literature)
% Details can be adjusted in macros.tex.
%
% BEWARE: Different fields of research and different university departments
% have their own customs regarding bibliography. Consult the bibliography
% format with your supervisor.
%
% The basic format according to the ISO 690 standard with numbered references
\usepackage[natbib,style=iso-numeric,sorting=none]{biblatex}
% ISO 690 with alphanumeric references (abbreviations of authors' names)
%\usepackage[natbib,style=iso-alphabetic]{biblatex}
% ISO 690 with references Author (year)
%\usepackage[natbib,style=iso-authoryear]{biblatex}
%
% Some fields of research prefer a simple format with numbered references
% (sorting=none tells that bibliography should be listed in citation order)
%\usepackage[natbib,style=numeric,sorting=none]{biblatex}
% Numbered references, but [1,2,3,4,5] is compressed to [1-5]
%\usepackage[natbib,style=numeric-comp,sorting=none]{biblatex}
% A simple format with alphanumeric references:
%\usepackage[natbib,style=alphabetic]{biblatex}

% Load the file with bibliography entries
\addbibresource{bibliography.bib}

% Our definitions of macros (see description inside)
\input macros.tex

%%% Title page and various mandatory informational pages
\begin{document}
%%% Title page of the thesis and other mandatory pages

%%% Inscriptions at the opening page of the hard cover

% We usually do not typeset the hard cover, but if you want to do it, change \iffalse to \iftrue
\iffalse

\pagestyle{empty}
\hypersetup{pageanchor=false}
\begin{center}

\large
Charles University

\medskip

Faculty of Mathematics and Physics

\vfill

{\huge\bf\ThesisTypeTitle}

\vfill

{\huge\bf\ThesisTitle\par}

\vfill
\vfill

\hbox to \hsize{\YearSubmitted\hfil \ThesisAuthor}

\end{center}

\newpage\openright
\setcounter{page}{1}

\fi

%%% Title page of the thesis

\pagestyle{empty}
\hypersetup{pageanchor=false}
\begin{center}

\centerline{\mbox{\includegraphics[width=166mm]{img/logo-en.pdf}}}

\vspace{-8mm}
\vfill

{\bf\Large\ThesisTypeTitle}

\vfill

{\LARGE\ThesisAuthor}

\vspace{15mm}

{\LARGE\bfseries\ThesisTitle\par}

\vfill

\Department

\vfill

{
\centerline{\vbox{\halign{\hbox to 0.45\hsize{\hfil #}&\hskip 0.5em\parbox[t]{0.45\hsize}{\raggedright #}\cr
Supervisor of the \ThesisTypeName{} thesis:&\Supervisor \cr
\ifx\ThesisType\TypeRig\else
\noalign{\vspace{2mm}}
Study programme:&\StudyProgramme \cr
\fi
}}}}

\vfill

Prague \YearSubmitted

\end{center}

\newpage

%%% A page with a solemn declaration to the thesis

\openright
\hypersetup{pageanchor=true}
\vglue 0pt plus 1fill

\noindent
I declare that I carried out this \ThesisTypeName{} thesis on my own, and only with the cited
sources, literature and other professional sources.
I understand that my work relates to the rights and obligations under the Act No.~121/2000 Sb.,
the Copyright Act, as amended, in particular the fact that the Charles
University has the right to conclude a license agreement on the use of this
work as a school work pursuant to Section 60 subsection 1 of the Copyright~Act.

\vspace{10mm}

\hbox{\hbox to 0.5\hsize{%
In \hbox to 6em{\dotfill} date \hbox to 6em{\dotfill}
\hss}\hbox to 0.5\hsize{\dotfill\quad}}
\smallskip
\hbox{\hbox to 0.5\hsize{}\hbox to 0.5\hsize{\hfil Author's signature\hfil}}

\vspace{20mm}
\newpage

%%% Dedication

\openright

\noindent
\Dedication

\newpage

%%% Mandatory information page of the thesis

\openright
{\InfoPageFont

\vtop to 0.5\vsize{
\setlength\parindent{0mm}
\setlength\parskip{5mm}

Title:
\ThesisTitle

Author:
\ThesisAuthor

\DeptType:
\Department

Supervisor:
\Supervisor, \SupervisorsDepartment

Abstract:
\Abstract

Keywords:
{\def\sep{\unskip, }\ThesisKeywords}

\vfil
}

% In Czech study programmes, it is mandatory to include Czech meta-data:

\ifx\StudyLanguage\LangCS

\vtop to 0.49\vsize{
\setlength\parindent{0mm}
\setlength\parskip{5mm}

Název práce:
\ThesisTitleCS

Autor:
\ThesisAuthor

\DeptTypeCS:
\DepartmentCS

Vedoucí bakalářské práce:
\Supervisor, \SupervisorsDepartmentCS

Abstrakt:
\AbstractCS

Klíčová slova:
{\def\sep{\unskip, }\ThesisKeywordsCS}

\vfil
}

\fi

}

\newpage

%%% Further pages will be numbered
\pagestyle{plain}


%%% A page with automatically generated table of contents of the thesis

\tableofcontents

%%% Each chapter is kept in a separate file
\chapter*{Introduction}
\addcontentsline{toc}{chapter}{Introduction}

Modern software systems are subject to continuous evolution driven by shifting user requirements. These changes often affect not only application logic but also the structure of the underlying data and the queries executed over it. Consequently, developers must adapt the data-access code in response to evolving schemas, altered query workloads, or even complete changes of the underlying database management system (DBMS). In the worst-case scenario, such changes require data migration across DBMS platforms and necessitate corresponding adaptations at the application level -- particularly in the mapping and querying logic embedded in object-relational mapping (ORM) frameworks.

While the DBMS side of the problem has received considerable attention -- with numerous academic and commercial solutions~\cite{TODO} for data migration, query rewriting, and query synchronization -- less effort has been directed toward managing the consequences of such changes on application code. Existing approaches~\cite{TODO} support, for example, the rewriting of SQL queries for target DBMSs or synchronization within a logical data model. However, adapting the application-side ORM layers (entity mappings, repository classes, query logic) remains a largely manual and error-prone process. This gap becomes especially critical when the original ORM framework, once optimal for a given workload, becomes inefficient or incompatible due to backend changes. In such cases, memory overhead, runtime performance, or even functional compatibility can be severely affected.

On the application side, the complexity is further amplified by the diversity of available ORM frameworks. Each framework offers different abstractions, feature sets, performance characteristics, and support for advanced mapping patterns. Developers typically commit to a single ORM framework, which can result in suboptimal performance if the selected framework does not align well with specific query workloads or with the target DBMS. In many scenarios, a hybrid approach -- utilizing different ORM frameworks for disjoint query subsets -- could yield significantly better runtime and memory efficiency.

Large language models (LLMs) have recently been explored~\cite{TODO} for translating between ORM frameworks. While promising, current LLM-based solutions remain limited in practice: translations may be incomplete, semantically incorrect, or based on deprecated or insecure libraries. More critically, LLMs do not support optimization objectives: they lack mechanisms to reason about execution cost or to recommend optimal target ORM configurations under resource constraints (e.g., maximum memory usage per query or across a query workload).

In this thesis... , we present \emph{ORMorpher}, a novel tool that supports both translation and optimization across ORM frameworks. Our contributions are as follows:
\begin{itemize}
    \item We introduce a translation framework for migrating entity and repository classes between ORM frameworks.
    \item We propose an optimization framework to automatically select the most efficient (combination of) ORM framework(s), or combination of frameworks, for a given set of queries and resource constraints.
    \item We demonstrate extensibility toward application-level migration strategies supporting DBMS transitions.
    \item We provide a prototype implementation targeting prominent .NET ORM frameworks (NHibernate, Entity Framework Core, and Dapper) over Microsoft SQL Server.
\end{itemize}

\paragraph{Outline} The rest of the thesis is organised as follows...

\chapter{ORM feature comparison}
\chapter{Experimental comparison}

This chapter presents an experimental comparison of the 7 selected ORM frameworks. 
The primary goals are to assert the correctness of results, measure performance, and assess how effectively each framework integrates into the .NET through its supported query language.

For each framework, we will create a set of unit tests performing pre-selected queries. They should cover different areas of querying through ORMs, such as entity mapping, relationships, result modification, and aggregation.

We will only test read (SELECT) queries. Those cover a sufficiently large area of features and provide us with enough information to get some insight on capabilities and performance. Other non-read operations, including insert, update, and delete queries, could be explored in future work.

\section{Dataset}\label{sec:dataset}
% https://learn.microsoft.com/en-us/sql/samples/wide-world-importers-what-is?view=sql-server-ver16
To adequately test performance, we need a sufficiently large data set. Microsoft offers pre-made open source datasets.
For this comparison, the Wide World Importers sample database \cite{microsoftWWI} has been selected.
It contains a diverse set of fictional data about a company Wide World Importers described as ``wholesale novelty goods importer and distributor operating from the San Francisco bay area.``\cite{microsoftWWI}

The database contains over thirty tables split into four schemas: Application, Purchasing, Sales, Warehouse. For each query, we will select a suitable table and columns that can best showcase the desired feature.

\section{Database}\label{sec:database}
% database, docker, test project, .net
The dataset is made for Microsoft SQL Server (MSSQL). It is the third most popular database management system as of March 2025 according to DB-Engines Ranking\cite{DBEngines2025Ranking}. With a score of 788.14, it is behind Oracle (1253.08) and MySQL (988.13). The score takes into account number of mentions on websites, frequency of technical discussions, number of job offers mentioning the system, relevance in social networks and other parameters.\cite{DBEnginesRankingMetrics} 
Microsoft SQL Server is commonly used in .NET applications with EF Core ORM as the three products are all developed by Microsoft.  

Unit tests should be easy to set up and run. Configuring a database can turn into a complicated task. We will create a Docker container running MSSQL. The dataset will be loaded on startup. Execution of unit tests will be just a matter of starting the Docker container and running the \texttt{dotnet test} command (or alternatively using Visual Studio Test Explorer). 

\section{Testing approach}\label{sec:testing_approach}
We will prefer to use LINQ where possible as it is tightly integrated into the .NET environment. 
If LINQ integration is not available, we will use any suitable alternative. And if there are no other possibilities, we will default to raw SQL queries. 

There are multiple different methods for mapping database tables to entities. We will be using different methods to show the variety and their impact or trade-offs. 

If ORM supports caching entities, we will disable it so that our benchmarks are not affected. 
A feature related to caching is change tracking. For example, EF Core tracks changes made to an entity, so that it can later perform an update. All our tests will only execute read operations. Change tracking would unnecessarily increase measured time or allocated memory, so we will disable it.

Logging will be turned off for performance tests. In feature tests, we might want to compare generated SQL. So, we will enable SQL logging into the test output window where possible.


\section{Selected queries}\label{sec:selected_queries}
Our goal is to test the query capabilities as broadly as possible. We want to create queries that use different conditions, result modifications, aggregations, relationships between tables, etc. The queries are grouped into categories A-H by the type of features they test.

Each query type will contain an example of what the query should look like in raw SQL. The example can be used to directly query the test database to preview results. The example might be simplified for this text. Unit tests might use a different query language where available.

\subsection{A - Entity projection}
This group will test how well ORM can handle projecting table columns onto a user-defined entity or entities.

\subsubsection*{A1 Entity identical to a table} \label{query:a1}
The test will retrieve a table row and map it to an entity with properties identical to table columns. 
Table \texttt{Purchasing.PurchaseOrders} will be queried and one item retrieved based on its ID.

\begin{lstlisting}[language=SQL]
SELECT * FROM Purchasing.PurchaseOrders 
WHERE PurchaseOrderId = 25
\end{lstlisting}

\subsubsection*{A2 Limited entity} \label{query:a2}
The resulting entity will have less properties than table columns. Only the data we really need should be transferred. 
Table \texttt{Purchasing.Suppliers} will be queried and only columns related to supplier's contact info will be retrieved.

\begin{lstlisting}[language=SQL]
SELECT SupplierID, SupplierName, PhoneNumber, FaxNumber, WebsiteURL, ValidFrom, ValidTo 
FROM Purchasing.Suppliers 
WHERE SupplierID = 10
\end{lstlisting}

\subsubsection*{A3 Multiple entities from one table} \label{query:a3}
One table will be queried and the result will be divided into two different entities. 
Table \texttt{Purchasing.Suppliers} will be used again, from which we will retrieve contact information and bank account information into two separate entities. 

\begin{lstlisting}[language=SQL]
SELECT 
    SupplierId, SupplierName, PhoneNumber, FaxNumber, WebsiteURL, ValidFrom, ValidTo, 
    SupplierId, BankAccountName, BankAccountBranch, BankAccountCode, BankAccountNumber, BankInternationalCode 
FROM Purchasing.Suppliers 
WHERE SupplierID = 10
\end{lstlisting}

\subsubsection*{A4 Stored procedure result into entity} \label{query:a4}
This query will execute a stored procedure, limited by parameters, and load the result into an entity.
The executed stored procedure will be \texttt{Integration.GetOrderUpdates} with parameters LastCutoff and NewCutoff. 
We will limit the cut-off to a one-year range from 2014 to 2015. 66741 order updates should be returned.

This stored procedure returns columns with spaces in their names, for example, ``WWI Order ID``. As properties in the C\# language cannot contain spaces, it will be interesting to see how and if different frameworks can handle this.

\begin{lstlisting}[language=SQL]
EXEC WideWorldImporters.Integration.GetOrderUpdates 
@LastCutoff = '2014-01-01', @NewCutoff = '2015-01-01'
\end{lstlisting}

% \paragraph{} % - jak tady vrátit tenhle text zpět 
\subsubsection*{A Summary}
Queries \textbf{A1}, \textbf{A2}, and \textbf{A3} will fetch one row based on its ID. The measured time and memory allocation will then show the overhead of the ORM framework when mapping data to the resulting entities. Query \textbf{A4} returns a large number of results, so it is a first query that shows us the performance with high-volume data.

\subsection{B - Selection}
Probably the most common query operation is limiting results based on a condition. This set of queries will query table \texttt{Sales.OrderLines} with varied conditions.

\subsubsection*{B1 Selection over indexed column} \label{query:b1}
The query will retrieve one order line based on the ID of its order. The ID is foreign key with an index built over it.

\begin{lstlisting}[language=SQL]
SELECT * FROM Sales.OrderLines WHERE OrderID = 26866
\end{lstlisting}

\subsubsection*{B2 Selection over non-indexed column} \label{query:b2}
Unlike the first query, this one will fetch order lines filtered by a column without an index. The column that we will filter over will be unit price of the order line.
It is worth noting that it should not be relevant if a column is indexed or not for ORM comparison. Because all the queries will be performed over the same database system. We include both queries in case an interesting result appears.

\begin{lstlisting}[language=SQL]
SELECT * FROM Sales.OrderLines WHERE UnitPrice = 25
\end{lstlisting}

\subsubsection*{B3 Range query} \label{query:b3}
The previous queries filtered based on a single value, but a range of values is often desirable as well.
This query will select a subset of order lines with a column \texttt{PickingCompletedWhen} within a selected date range.

\begin{lstlisting}[language=SQL]
SELECT * FROM Sales.OrderLines 
WHERE PickingCompletedWhen 
BETWEEN '2014-12-20' AND '2014-12-31'
\end{lstlisting}

\subsubsection*{B4 In query} \label{query:b4}
When a range is not sufficient and we want specific values not fitting into an interval, we can provide a collection of values.
This query will return order lines belonging to orders with IDs 1, 10, 100, 1000, and 10000.

\begin{lstlisting}[language=SQL]
SELECT * FROM Sales.OrderLines 
WHERE OrderID IN (1, 10, 100, 1000, 10000)
\end{lstlisting}

\subsubsection*{B5 Text search} \label{query:b5}
This query will test a text search, looking for any order lines with description containing the word ``C++``.  The column does not have a full text search or any other index.

\begin{lstlisting}[language=SQL]
SELECT * FROM Sales.OrderLines 
WHERE Description LIKE '%C++%'
\end{lstlisting}

\subsubsection*{B6 Paging} \label{query:b6}
When filtering with a broad condition, a large amount of data might be returned. So we may opt to receive the data in batches. Another use case might be a paged table displayed to a user. For these a paging query with skip and take might be useful.
This query will fetch all order lines ordered by ID, but will skip first 1000 and take only next 50.

\begin{lstlisting}[language=SQL]
SELECT * FROM Sales.OrderLines 
ORDER BY OrderLineID 
OFFSET 1000 ROWS FETCH NEXT 50 ROWS ONLY
\end{lstlisting}

\subsection{C - Aggregation} \label{query:}
Aggregation functions are often used when performing calculations on data. We will test a select few.

\subsubsection*{C1 Count} \label{query:c1}
Count is very useful when we need to know the number of records that correspond to a specific condition. We will combine it with \texttt{GROUP BY} clause to retrieve distinct tax rates and the number of their occurrence.

\begin{lstlisting}[language=SQL]
SELECT TaxRate, COUNT(TaxRate) as Count 
FROM Sales.OrderLines 
GROUP BY TaxRate 
ORDER BY Count DESC
\end{lstlisting}

\subsubsection*{C2 Max}  \label{query:c2}
Maximum and minimum are also commonly used. We will attempt to retrieve the maximum unit price for all order lines.

\begin{lstlisting}[language=SQL]
SELECT MAX(UnitPrice) FROM Sales.OrderLines
\end{lstlisting}

\subsubsection*{C3 Sum} \label{query:c3}
Aggregation functions also allow performing operations. In this test, we will combine the \texttt{SUM} function with a multiplication inside. The result will be the total price across all order lines.
\begin{lstlisting}[language=SQL]
SELECT SUM(Quantity * UnitPrice) FROM Sales.OrderLines
\end{lstlisting}

\subsection{D - Relations}
An important feature of ORM should be mapping relations. Our domain entities are often connected, representing relations in real life.
However, as we found out in our previous analysis, not all the selected frameworks support mapping relations. 

Our test data contain no one to one relation. There should be a small difference between mapping one to one and one to many. We will cover the other types in sufficient detail.

Some ORMs support lazy loading. That means loading entities in a relation only after accessing the corresponding collection. That would result in more queries being sent and a potential slowdown.
We will prevent this by setting the loading strategy to eager when possible.

\subsubsection*{D1 One to many} \label{query:d1}
In one to many relation we have one parent and multiple children. A parent entity should have a collection containing its children. Optionally, a child entity can have a reference to its parent.
In this test, we will use the relation between order and its order lines. We will fetch one specific order along with all its lines.

We need to verify all child entities were returned and the parent collection filled. 



\begin{lstlisting}[language=SQL]
SELECT o.*, ol.* FROM Sales.Orders o
LEFT JOIN Sales.OrderLines ol ON ol.OrderID = o.OrderID
WHERE o.OrderID = 530
\end{lstlisting}

\subsubsection*{D2 Many to many} \label{query:d2}
In many to many relation, both sides can be connected with multiple entities of the same type. Therefore, it will be represented by a collection on both sides. The collection can contain the entity directly, which makes it easier for developers to work with. If necessary, it can contain entities that represent the join table for this relationship.

This test will consist of two queries. We want to test accessing the relationship from both sides. The first query will fetch stock items with their stock groups. And the other one stock groups with their stock items.
The relationship is represented by a join table \texttt{Warehouse.StockItemStockGroups}. If possible, we will try to map this relation without a join entity in our tests.

\begin{lstlisting}[language=SQL]
SELECT si.*, sg.* FROM Warehouse.StockItems si
LEFT JOIN Warehouse.StockItemStockGroups sisg
    ON si.StockItemID = sisg.StockItemID
LEFT JOIN Warehouse.StockGroups sg
    ON sisg.StockGroupID = sg.StockGroupID
ORDER BY si.StockItemID
\end{lstlisting}
\begin{lstlisting}[language=SQL]
 SELECT sg.*, si.* FROM Warehouse.StockGroups sg
 LEFT JOIN Warehouse.StockItemStockGroups sisg
     ON sg.StockGroupID = sisg.StockGroupID
 LEFT JOIN Warehouse.StockItems si
     ON sisg.StockItemID = si.StockItemID
 ORDER BY sg.StockGroupID
\end{lstlisting}
The queries utilise left join, so we will get all stock items, even if they belong to no stock groups. And vice versa for the other query.

\subsubsection*{D3 One to many with optional relation} \label{query:d3}
In the D1 test case, the one to many relation was not optional. Each order has at least one order line. In the D2 test case, each item belonged to at least one category as well.
This test case will utilize an optional one to many relation between a customer and a transaction. 

\begin{lstlisting}[language=SQL]
SELECT c.*, ct.* FROM Sales.Customers c
LEFT JOIN Sales.CustomerTransactions ct
    ON c.CustomerID = ct.CustomerID
ORDER BY c.CustomerID
\end{lstlisting}

\subsection{E - Result modification}
This category will test modifying the result set. We will test sorting by a column and filtering distinct result values.

\subsubsection*{E1 Column sorting} \label{query:e1}
Although we have used sorting to order the results of some previous tests to get deterministic results, this test will focus solely on sorting.
We will sort all purchase orders by their expected delivery date and retrieve only the first one thousand.

\begin{lstlisting}[language=SQL]
SELECT TOP (1000) * FROM Purchasing.PurchaseOrders 
ORDER BY ExpectedDeliveryDate ASC
\end{lstlisting}

\subsubsection*{E2 Distinct results} \label{query:e2}
To test fetching only unique values, we will request all supplier references from purchase orders.

\begin{lstlisting}[language=SQL]
SELECT DISTINCT SupplierReference 
FROM Purchasing.PurchaseOrders
\end{lstlisting}

\subsection{F - Querying JSON}
Microsoft SQL Server supports storing and querying JSON data\cite{mssqljson}. It is worth testing if our selected ORMs support constructing queries into the stored JSON.
Table \texttt{Application.People} contains text columns storing JSON values.

During the static comparison, we found that none of the selected ORMs supports XML or any other data format. So we will not be testing for those.

\subsubsection*{F1 JSON object query} \label{query:f1}

Column \texttt{CustomFields} contains a JSON with many different values. In this test we will focus on property \texttt{Title}. We will query all people with a title set to ``Team Member``.
We will accomplish that using the \texttt{JSON\_VALUE} SQL function.

\begin{lstlisting}[language=SQL]
SELECT * FROM Application.People
WHERE JSON_VALUE(CustomFields, '$.Title') = 'Team Member'
ORDER BY PersonId
\end{lstlisting}

\subsubsection*{F2 JSON array query} \label{query:f2}

The first case covered querying a JSON object. This one will test querying inside a JSON array. Column \texttt{OtherLanguages} contains a JSON array containing language names.

\begin{lstlisting}[language=SQL]
SELECT * FROM Application.People
WHERE EXISTS (
    SELECT 1
    FROM OPENJSON(OtherLanguages)
    WHERE value = 'Slovak'
)
\end{lstlisting}

\subsection{G - Set operations}

In the last group we will test common set operations. Those are useful when having to combine two sets of data.

\subsubsection*{G1 Union} \label{query:g1}

Union will combine results from both queries. In this test, we will combine results from two different intervals of IDs and ensure that they are all present in the result.

\begin{lstlisting}[language=SQL]
SELECT SupplierID FROM Purchasing.Suppliers 
    WHERE SupplierID < 5
UNION
SELECT SupplierID FROM Purchasing.Suppliers 
    WHERE SupplierID BETWEEN 5 AND 10
ORDER BY SupplierID
\end{lstlisting}

\subsubsection*{G2 Intersection} \label{query:g2}

Intersection will include results that appear in both of the queries. We will try combining supplier IDs in two overlapping intervals and assert that only the overlap is included. 

\begin{lstlisting}[language=SQL]
SELECT SupplierID FROM Purchasing.Suppliers 
    WHERE SupplierID < 10
INTERSECT
SELECT SupplierID FROM Purchasing.Suppliers 
    WHERE SupplierID BETWEEN 5 AND 15
ORDER BY SupplierID
\end{lstlisting}

\subsection{H - Querying metadata}
In the last section, we will briefly test the ORM framework's ability to query information about the database itself. While not commonly used in applications, such data can be required in specific cases.

\subsubsection*{H1 Column data type} \label{query:h1}
The final query will attempt to retrieve the data type of a table column. All information about the database schema is available through \texttt{INFORMATION\_SCHEMA.COLUMNS} table. The selected column has \texttt{NVARCHAR} type. The table contains more information like nullability, maximum length, etc. But for demonstration purposes, this query should be sufficient.

\begin{lstlisting}[language=SQL]
SELECT DATA_TYPE FROM INFORMATION_SCHEMA.COLUMNS 
WHERE TABLE_SCHEMA = 'Purchasing'
    AND TABLE_NAME = 'Suppliers'
    AND COLUMN_NAME = 'SupplierReference'
\end{lstlisting}

\section{Summary of the selected queries}
We tried to pick as broad a selection of features as possible. 
We certainly did not cover all of them, but our set should be representative enough to give us a fair idea of how different ORMs perform. 

The tests perform queries across different tables and fetch different amounts of data. 
Measured time and memory allocation will not allow us to make any assumptions when comparing different test cases for the same ORM.
But we will be able to compare identical test cases across different ORMs.

We can expect to see a difference in time performance between micro frameworks made to be efficient and feature-rich macro frameworks.
The micro frameworks utilize raw SQL and have very little overhead. When using a macro framework, the LINQ query has to be compiled into SQL. Rich mapping capabilities may also cause significant slowdown. And extra features may even cause higher memory consumption.

\section{Implementation}
The comparison was implemented in the form of a Visual Studio solution, containing a project folder for each framework. Each project folder contains three projects - entities, feature tests, and performance benchmark. 


\usetikzlibrary{positioning, arrows}
\begin{center}
\begin{tikzpicture}[
    node distance=1cm and 2cm,
    every node/.style={draw, minimum width=1cm, minimum height=1cm, align=center},
    every path/.style={->, >=stealth},
    faded/.style={draw=gray, fill=gray!20},
    faded_arrow/.style={->, >=stealth, draw=gray}
]

  \node (E) {Entities};
  \node (C) [right=of E, faded] {Common};
  \node (FT) [below left=of E] {Feature tests};
  \node (PB) [below right=of E] {Performance benchmark};

  \draw[thick, ->, >=stealth] (FT) -- (E);
  \draw[thick, ->, >=stealth] (PB) -- (E);
  \draw[faded_arrow] (FT) -- (C);
  \draw[faded_arrow] (PB) -- (C);

\end{tikzpicture}
\end{center}

Common project is also included. It contains a single class handling loading the database connection string from a configuration file. There is no dependency between feature tests and performance tests. It could be argued that the code for queries could be extracted to a separate class to avoid duplication. But we have kept it duplicated to isolate unit tests from performance benchmarks. This structure is repeated for each of the seven frameworks we are examining. There is again no dependency between the different frameworks for isolation purposes.

\subsection{Technology overview}
As we have decided before implementation in \autoref{sec:database}, we are using Microsoft SQL Server database with World Wide Importers data set. The database is running in a Docker container. Base image used is \texttt{mssql/server:2022-latest}\cite{mssqlDocker} provided by Microsoft. The image imports our dataset and waits until it is ready.

The .NET tests are structured in .NET projects grouped in one solution. .NET version 8 was chosen as it is the latest LTS (long term support) version at the time of writing.\cite{netversions} Tests were developed and run using Microsoft Visual Studio 2022, but \texttt{dotnet} CLI commands will be provided as well.

\subsection{Installation instructions}
The container was tested both in Docker and Podman. A Dockerfile is provided in benchmarks/Dockerfile. Container image can be built and started using the two commands below. Replace `docker` with `podman` if needed. The database will be exposed on port 1444, under username SA and password \texttt{Testingorms123}.
The third command can be used to test connection. The sqlcmd utility is available only if you have MSSQL installed locally. Otherwise, use any database management tool of your choice.

\begin{lstlisting}[language=sh]
docker build -t orm-comparison .

docker run -d --name orm-comparison -p 1444:1433 orm-comparison

sqlcmd -S 127.0.0.1,1444 -U SA -P Testingorms123 -Q "SELECT * FROM [WideWorldImporters].[Purchasing].[PurchaseOrders]"
\end{lstlisting}

To run feature unit tests, build the project in Visual Studio and then use the Test Explorer window to select and run specific tests. Alternatively, navigate to the folder with solution file and run \lstinline{dotnet test}. This will build the project, find all tests, and run them. 

To run performance benchmarks, set \texttt{BenchmarkMain} as a startup project. Right-click it in the Solution Explorer window and choose `Set as Startup project`. Make sure configuration is set to Release, not Debug. Then start the project without debugging (CTRL + F5 keys). A console window will appear. It contains instructions on how to target specific benchmarks. To run all, type in asterisk (*) and press enter. This will trigger full run, which takes approximately one hour to finish. You can switch to test configuration which does only a few number of iterations in \texttt{BenchmarkMain/Program.cs} file. What the configuration does is further explained in following sections. 

Console command for running benchmarks is \lstinline{dotnet run --configuration Release --project BenchmarkMain\BenchmarkMain.csproj}. To run with shorter test configuration append \lstinline{-- --testb} at the end of the command.


\subsection{Entities}
Set of classes either mirroring a database table or containing a necessary subset of properties. The columns mapped always depend on the test they are being used in. The entities are POCOs (plain old CLR objects). The term is derived from POJO (plain old Java objects). \cite{Fowler2003POJO} It refers to classes that don't inherit from any base class or interface. And they are not tied to any specific framework.

Below is an example of a POCO class for a purchase order. Visibility modifiers and some properties were omitted for readability. 

% TODO Move explanation about what entities are to earlier section/introduction

\begin{lstlisting}[language=CSharp]
class PurchaseOrder
{
    int PurchaseOrderID { get; set; }
    int SupplierID { get; set; }
    DateTime OrderDate { get; set; }
    DateTime? ExpectedDeliveryDate { get; set; }
    string? Comments { get; set; }
}
\end{lstlisting}

Some frameworks like Dapper do not need any entity mapping and simply require the properties to match selected query columns. The purchase order entity could be used to load results of a query, but the query has to return columns with matching names, data types, and nullability.

More complex frameworks that do not just perform read queries need to know how to map an entity to a database table. The corresponding table might not even exist and is generated based on the entity. Such use case is central to domain-driven development\cite{FowlerDDD}.

Entity mapping might be done in multiple ways. The simplest form is no mapping. The framework would then guess table and column names, data types, nullability, etc. from the way an entity is named and declared. Clearly, that lacks details like numeric precision, constraints, and indices.
Those details could be specified by property and class attributes. That makes them tied to an entity, which could be a problem if we want to isolate the domain model from database access. 
Mapping through code or a configuration file allows separation and can be more expressive and powerful.

To showcase different approaches to mapping, we have tried to select a different method for each framework test. Mapping is usually compiled on application startup or first query and then cached. So performance should not be affected by a choice of mapping. Dapper does not use any mapping (as it does not support any). PetaPoco and RepoDB use attributes to declare table schema and name, and in some cases column names, the rest is inferred. For those two, mapping abilities are quite limited. LINQ to DB utilizes mapping by code through builder pattern. NHibernate has its proprietary mapping through files in XML format. Both versions of Entity Framework use attribute mapping combined with mapping through code. Their attribute mapping is much more powerful than those of PetaPoco and RepoDB, but in more complex cases, mapping by code has to be supplied.

First listing below shows attribute mapping in EF Core, where attributes are placed directly above classes and properties. In the second one, part of an XML mapping file for the same entity is shown. The verbosity of the mapping is noticeable.

\begin{lstlisting}[language=CSharp,basicstyle=\ttfamily\footnotesize]
[Table("PurchaseOrders", Schema = "Purchasing")]
class PurchaseOrder
{
    [Key]
    int PurchaseOrderID { get; set; }
}
\end{lstlisting}

\begin{lstlisting}[language=xml,basicstyle=\ttfamily\footnotesize]
<?xml version="1.0" encoding="utf-8" ?>
<hibernate-mapping xmlns="urn:nhibernate-mapping-2.2" namespace="NHibernateEntities">
    <class name="NHibernateEntities.PurchaseOrder, NHibernateEntities" table="PurchaseOrders" schema="Purchasing">
        <id name="PurchaseOrderID" column="PurchaseOrderID" type="int">
            <generator class="identity" />
        </id>
        <property name="SupplierID" not-null="true" />
        <property name="OrderDate" not-null="true" />
        ...
    </class>
</hibernate-mapping>
\end{lstlisting}

\subsection{Feature tests}

For each query defined in \autoref{sec:selected_queries}, a unit test in the form of an isolated method was implemented. Connection and other configuration are created outside of the test in a setup part. The method executes the query and correct results are asserted. Assertion checks the amount of results and their order. Each property is tested to contain correct data to ensure mapping was done correctly.

The example below shows Query A1 (\autoref{query:a1}) implemented for Dapper. A single purchase order is retrieved based on its ID, and its properties are checked for correct values. 
\begin{lstlisting}[language=CSharp]
[Fact]
public void A1_EntityIdenticalToTable()
{
    var order = connection.QuerySingle<PurchaseOrder>(
        "SELECT * FROM WideWorldImporters.Purchasing.PurchaseOrders WHERE PurchaseOrderId = @PurchaseOrderId",
        new { PurchaseOrderId = 25 }
    );

    Assert.Multiple(() =>
    {
        Assert.Equal(25, order.PurchaseOrderID);
        Assert.Equal(12, order.SupplierID);
        Assert.Equal(new DateTime(2013, 1, 5), order.OrderDate);
        ...
    });
}
\end{lstlisting}

As decided in \autoref{sec:testing_approach} we have attempted to use LINQ query language. Where not supported we have used any available alternative that did not involve just writing raw SQL. And finally, where nothing else was possible we have used the raw alternative to have a passing test.  
%TODO mention the table that will have feature comparison

\subsection{Performance benchmarks}
Feature tests can assert the correctness of returned results and overall functionality. But their execution time can greatly vary. There can be some overhead, coming from test initialization or assertions. The system the benchmarks are running on can also impact results.

To solve these problems we have used BenchmarkDotNet library\cite{BenchmarkDotNet}. It allows writing benchmark methods very similar to unit tests. In fact, they are completely identical to unit tests, the only differences are missing assertions and a different method attribute. It allows us to ``transform methods into benchmarks, track their performance, and share reproducible measurement experiments`` \cite{BenchmarkDotNet}. Work with the library was quite easy, it is very configurable and easy to run. 

It runs benchmark for each method multiple times. It starts with a few iterations which will not be included in the results. The first few ensure just-in-time compilation has taken effect, then multiple iterations without actually running the method code are started to compile overhead. The overhead is removed from the results, ensuring reliability. 


According to BenchmarkDotNet documentation\cite{BenchmarkDotNetHow} a single benchmark run consists of the following stages.
\begin{itemize}
    \item OverheadJitting - Measures benchmarking infrastructure compilation.
    \item WorkloadJitting - Measures benchmark method compilation.
    \item WorkloadPilot - Selects best iteration count based on heuristics.
    \item WorkloadWarmup - Runs several iterations before measuring to ensure factors like CPU caching or branch prediction are optimized.
    \item WorkloadActual - Actual measurements
    \item WorkloadResult - Result calculated from the actual run without the overhead.
\end{itemize}

After the benchmark run finishes, results and run log can be found in:

\smallskip
\noindent\path{benchmarks/BenchmarkMain/bin/Release/net8.0/BenchmarkDotNet.Artifacts}

\smallskip
The referenced run results are copied to:

\smallskip
\noindent\path{benchmarks/results/joined}

\smallskip
Notable result files include:
\begin{itemize}
    \item \path{BenchmarkRun-joined-2025-03-06-16-10-18-report.csv}\\
    Measurements of all benchmarks in CSV format.

    \item \path{BenchmarkRun-joined-2025-03-06-16-10-18-report.html}\\
    Summary table of all benchmarks (HTML format).

    \item \path{BenchmarkRun-joined-2025-03-06-16-10-18-report-github.md}\\
    Summary table of all benchmarks (Markdown format suitable for GitHub).

    \item \path{BenchmarkRun-joined-2025-03-06-16-10-18-measurements.csv}\\
    Detailed measurements of every single iteration of each benchmark.
\end{itemize}

\section{Results}
Results are accompanied by tables showing results across all test cases. 

\autoref{tab:feature_comp} showcases feature test results. The base part of the result shows if the query could be expressed in the framework, Y means yes and N means no. Further detail is then clarified. The categories for successful results are, sorted from the most ideal to the least:
\begin{itemize}
    \item \textbf{Y LINQ} - The query could be written in LINQ with no issue.
    \item \textbf{Y lambda expr.} - The query could be written with a lambda expression. Similar to LINQ but can express less.
    \item \textbf{Y raw SQL} - No statically typed language feature could be used and the query had to be written in SQL. Either in string or using SQL builder.
    \item \textbf{N manually in memory} - A different or limited query had to be used and the result had to be modified or computed in memory.
    \item \textbf{N custom JSON parsing/converter} - JSON could not be queried directly.
    \item \textbf{N error} - The query threw an exception and it was not possible to fix. 
\end{itemize}
The negative results are further described in the text.

\autoref{tab:benchmark_results_time} shows measured benchmark time in microseconds (one millionth of a second). And \autoref{tab:benchmark_results_memory} shows allocated memory during the run of the benchmark method. The results presented in text will be rounded to the nearest integer.

Performance tests have to be considered in the context of the feature test results. It was not possible to always write the test in the same way across all frameworks. In a few cases, it was not possible to do the query using the framework at all. In those cases, to have a passing test, pure ADO.NET query was used. 

In some cases it was not possible to get the desired result purely as a result of a query. Some relations or calculations had to be done in memory. That, of course, majorly affects measured memory consumption. Such anomalies will always be described in the corresponding test results.

Certain frameworks do not allow logging of generated SQL statements to console or file. We can work around this using SQL Server Profiler\cite{sqlProfiler}. It is a tool integrated into SQL Server Management Studio. It allows for the capturing of any event happening in the SQL Server database, including incoming SQL queries. That makes us able to inspect the raw query after it leaves our application and the frameworks can make no further edit to it.

\subsection{Benchmark execution}
The final performance benchmark run was executed on a PC with the highest performance power plan while being constantly plugged in. Screen saving was turned off. The only applications running were those essential to the benchmark (database in Docker container and Visual Studio project). No other work was being done. Data further present are a result of that run and are available as an attachment.

The exact PC and runtime specifications are available below. This information is comes directly from the benchmark run log. 
\begin{lstlisting}
BenchmarkDotNet v0.14.0, Windows 11 (10.0.22631.4890/23H2/2023Update/SunValley3)
AMD Ryzen 5 5600H with Radeon Graphics, 1 CPU, 12 logical and 6 physical cores
.NET SDK 9.0.101
  [Host]     : .NET 8.0.11 (8.0.1124.51707), X64 RyuJIT AVX2
  DefaultJob : .NET 8.0.11 (8.0.1124.51707), X64 RyuJIT AVX2
\end{lstlisting}

\subsection{A - Entity projection}
The first three queries \hyperref[query:a1]{A1}, \hyperref[query:a2]{A2} and \hyperref[query:a3]{A3} fetched one database record and shaped it in different ways. All ORMs were capable of this task without any issue. It was a matter or having an entity with the chosen columns and shaping a query. The macro frameworks like EF and NHibernate allowed us to set the shape using \texttt{Select} LINQ method. In the micro ones, we had to choose the columns through \texttt{SELECT} part of the SQL statement. 
There are no major differences in the measured time, just a slight increase for the macro ORMs as would be expected. When looking at allocated memory, we can see a 2-10$\times$ increase for NHibernate  and especially EF6. EF Core is clearly optimized in this area and the increase in allocated memory is not major. 

We encounter a first issue at the \hyperref[query:a4]{A4} query which loads the result of a stored procedure. The stored procedure in our test data returns columns with spaces in their names. While it's not a usual naming convention, it was not expected to make some frameworks completely fail.

Dapper, PetaPoco, RepoDB, linq2db and EF Core had no issue. All of those had to use a raw SQL to specify the query name, there was no way to map it to for example an object. For the context, over 60 000 results were returned. Measured time has no significant difference, but allocated memory differs by a huge margin. PetaPoco, RepoDB and linq2db allocated 16-17 MB, while Dapper and EF Core allocated 24 and 29 MB respectively. After checking with SQL Server Profiler\cite{sqlProfiler}, there is no major difference in the executed query. It always returns the same columns. And our assertions made sure all properties are filled in.

NHibernate was failing on an internal error when converting the stored procedure result.
The error was \lstinline{NHibernate.PropertyNotFoundException : Could not find a setter for property 'WWI Order ID' in class 'System.RuntimeType'}.
It suggests its internal representation of the entity could not be filled due to the white space. The only fix we could find was converting the result into a hash table instead of a list of entities. As seen in the measurements, it increased both time and allocated memory significantly compared to others.

And finally, EF6 had the same issue with spaces, except we could not find any workaround in this case. There is no similar feature to load the result into a hash table, or dynamic objects. So the test uses a pure ADO.NET connection, with no EF6 features whatsoever. We consider the test failed and the measurements are not relevant for our comparison.

\subsection{B - Selection}
This group tests different types of conditions. Feature-wise, all tests pass. There is no exception or difference. As expected, queries in Dapper and PetaPoco have to be written in SQL. While in the other frameworks, we can utilize LINQ (or lambda expressions) to have statically typed queries. 

While the tests do pass, measurements differ a lot per ORM and query. Let us examine the queries one by one.

As we have discussed when defining the queries, it is not relevant for ORM comparison using the same relational database if a column is indexed or not. We will therefore not compare query B1 and B2 in this sense. However, the first query only returns 5 rows, while the second over thirteen thousand.

Query \hyperref[query:b1]{B1} has relatively little difference in terms of measured time. The macro ORMs again show a small increase. But the big difference is in the memory consumption. EF6 allocated 102 KB and NHibernate 49 KB. The other frameworks allocated in the range of 8 to 15 KB.
Especially interesting is the memory consumption between Dapper (9 KB) and PetaPoco (14 KB). While taking almost the same time, PetaPoco allocated two times more memory. 
% TODO profile memory

The second query \hyperref[query:b2]{B2} shows unusual results in measured time, NHibernate took 2 times more and EF6 3 times more time than other frameworks. After checking with SQL profiler, there is no difference in SQL queries that would affect performance or the amount of data returned.

In terms of memory, we can again see the spike for NHibernate and EF6, allocating 7 and 16 times more respectively. If we compare Dapper and PetaPoco, this time Dapper allocated approximately two times more memory. 
This might suggest PetaPoco has a bigger framework overhead, which shows at small results. But for bigger amounts of objects, it has better memory optimization. 

The third query \hyperref[query:b3]{B3} took about 22 milliseconds for most frameworks. Dapper, PetaPoco and RepoDB, the three smaller ORMs, took about 30 milliseconds. There is an interesting difference in the query. In Dapper and PetaPoco, where SQL is written by us, we used \texttt{BETWEEN} condition for selecting a date range. Other ORMs broke this condition down into separate less than or equal to and greater than or equal to parts. This could hint at what the reason for slower execution might be. But RepoDB does not use \texttt{BETWEEN} making this assumption invalid. Profile tool shows no other difference in the generated query. As for memory consumption, we can see significant spike for NHibernate and EF6.

The query \hyperref[query:b4]{B4} fetched 15 results from a table with over 231 000 rows using the IN statement on primary key.



% possible TODO find sources comparing SELECT * vs listing all columns 


%\clearpage
\afterpage{
\begin{landscape}
%If the table is too wide, replace \begin{table}[!htp]...\end{table} with
%\begin{adjustwidth}{-2.5 cm}{-2.5 cm}\centering\begin{threeparttable}[!htb]...\end{threeparttable}\end{adjustwidth}
%\begin{adjustwidth}{-2.5 cm}{-2.5 cm}
\begin{table}[htp]
\centering
\caption{Feature comparison}
\label{tab:feature_comp}
\scriptsize
\begin{threeparttable}[!htb]
\def\arraystretch{1.25}
\begin{tabular}{
>{\raggedright\arraybackslash}p{50.00mm}
>{\raggedright\arraybackslash}p{20.00mm}
>{\raggedright\arraybackslash}p{20.00mm}
>{\raggedright\arraybackslash}p{20.00mm}
>{\raggedright\arraybackslash}p{20.00mm}
>{\raggedright\arraybackslash}p{20.00mm}
>{\raggedright\arraybackslash}p{20.00mm}
>{\raggedright\arraybackslash}p{20.00mm}
>{\raggedright\arraybackslash}p{20.00mm}
}
\toprule
\textbf{Benchmark Name} & \textbf{Dapper} & \textbf{PetaPoco} & \textbf{RepoDB} & \textbf{linq2db} & \textbf{NHibernate} & \textbf{EF6} & \textbf{EF Core} \\
\midrule
A1\_EntityIdenticalToTable &Y raw SQL &Y raw SQL &\cellcolor[HTML]{d9ead3}Y lambda expr. &\cellcolor[HTML]{b7e1cd}Y LINQ &\cellcolor[HTML]{b7e1cd}Y LINQ &\cellcolor[HTML]{b7e1cd}Y LINQ &\cellcolor[HTML]{b7e1cd}Y LINQ \\
A2\_LimitedEntity &Y raw SQL &Y raw SQL &\cellcolor[HTML]{d9ead3}Y lambda expr. &\cellcolor[HTML]{b7e1cd}Y LINQ &\cellcolor[HTML]{b7e1cd}Y LINQ &\cellcolor[HTML]{b7e1cd}Y LINQ &\cellcolor[HTML]{b7e1cd}Y LINQ \\
A3\_MultipleEntitiesFromOneResult &Y raw SQL &Y raw SQL &\cellcolor[HTML]{d9ead3}Y lambda expr. &\cellcolor[HTML]{b7e1cd}Y LINQ &\cellcolor[HTML]{b7e1cd}Y LINQ &\cellcolor[HTML]{b7e1cd}Y LINQ &\cellcolor[HTML]{b7e1cd}Y LINQ \\
A4\_StoredProcedureToEntity &Y raw SQL &Y raw SQL &Y raw SQL &Y raw SQL &\cellcolor[HTML]{ea9999}N error &\cellcolor[HTML]{ea9999}N error &Y raw SQL \\
B1\_SelectionOverIndexedColumn &Y raw SQL &Y raw SQL &\cellcolor[HTML]{d9ead3}Y lambda expr. &\cellcolor[HTML]{b7e1cd}Y LINQ &\cellcolor[HTML]{b7e1cd}Y LINQ &\cellcolor[HTML]{b7e1cd}Y LINQ &\cellcolor[HTML]{b7e1cd}Y LINQ \\
B2\_SelectionOverNonIndexedColumn &Y raw SQL &Y raw SQL &\cellcolor[HTML]{d9ead3}Y lambda expr. &\cellcolor[HTML]{b7e1cd}Y LINQ &\cellcolor[HTML]{b7e1cd}Y LINQ &\cellcolor[HTML]{b7e1cd}Y LINQ &\cellcolor[HTML]{b7e1cd}Y LINQ \\
B3\_RangeQuery &Y raw SQL &Y raw SQL &\cellcolor[HTML]{d9ead3}Y lambda expr. &\cellcolor[HTML]{b7e1cd}Y LINQ &\cellcolor[HTML]{b7e1cd}Y LINQ &\cellcolor[HTML]{b7e1cd}Y LINQ &\cellcolor[HTML]{b7e1cd}Y LINQ \\
B4\_InQuery &Y raw SQL &Y raw SQL &\cellcolor[HTML]{d9ead3}Y lambda expr. &\cellcolor[HTML]{b7e1cd}Y LINQ &\cellcolor[HTML]{b7e1cd}Y LINQ &\cellcolor[HTML]{b7e1cd}Y LINQ &\cellcolor[HTML]{b7e1cd}Y LINQ \\
B5\_TextSearch &Y raw SQL &Y raw SQL &\cellcolor[HTML]{d9ead3}Y lambda expr. &\cellcolor[HTML]{b7e1cd}Y LINQ &\cellcolor[HTML]{b7e1cd}Y LINQ &\cellcolor[HTML]{b7e1cd}Y LINQ &\cellcolor[HTML]{b7e1cd}Y LINQ \\
B6\_PagingQuery &Y raw SQL &Y raw SQL &\cellcolor[HTML]{d9ead3}Y lambda expr. &\cellcolor[HTML]{b7e1cd}Y LINQ &\cellcolor[HTML]{b7e1cd}Y LINQ &\cellcolor[HTML]{b7e1cd}Y LINQ &\cellcolor[HTML]{b7e1cd}Y LINQ \\
C1\_AggregationCount &Y raw SQL &Y raw SQL &\cellcolor[HTML]{f4cccc}N manually in memory &\cellcolor[HTML]{b7e1cd}Y LINQ &\cellcolor[HTML]{b7e1cd}Y LINQ &\cellcolor[HTML]{b7e1cd}Y LINQ &\cellcolor[HTML]{b7e1cd}Y LINQ \\
C2\_AggregationMax &Y raw SQL &Y raw SQL &\cellcolor[HTML]{d9ead3}Y lambda expr. &\cellcolor[HTML]{b7e1cd}Y LINQ &\cellcolor[HTML]{b7e1cd}Y LINQ &\cellcolor[HTML]{b7e1cd}Y LINQ &\cellcolor[HTML]{b7e1cd}Y LINQ \\
C3\_AggregationSum &Y raw SQL &Y raw SQL &Y raw SQL &\cellcolor[HTML]{b7e1cd}Y LINQ &\cellcolor[HTML]{b7e1cd}Y LINQ &\cellcolor[HTML]{b7e1cd}Y LINQ &\cellcolor[HTML]{b7e1cd}Y LINQ \\
D1\_OneToManyRelationship &\cellcolor[HTML]{f4cccc}N manually in memory &\cellcolor[HTML]{f4cccc}N manually in memory &\cellcolor[HTML]{d9ead3}Y lambda expr. &\cellcolor[HTML]{b7e1cd}Y LINQ + mapping &\cellcolor[HTML]{b7e1cd}Y LINQ + mapping &\cellcolor[HTML]{b7e1cd}Y LINQ + mapping &\cellcolor[HTML]{b7e1cd}Y LINQ + mapping \\
D2\_ManyToManyRelationship &\cellcolor[HTML]{f4cccc}N manually in memory &\cellcolor[HTML]{f4cccc}N manually in memory &\cellcolor[HTML]{f4cccc}N manually in memory &\cellcolor[HTML]{b7e1cd}Y LINQ + mapping, join entity &\cellcolor[HTML]{b7e1cd}Y LINQ + mapping &\cellcolor[HTML]{b7e1cd}Y LINQ + mapping &\cellcolor[HTML]{b7e1cd}Y LINQ + mapping \\
D3\_OptionalRelationship &\cellcolor[HTML]{f4cccc}N manually in memory &\cellcolor[HTML]{f4cccc}N manually in memory &\cellcolor[HTML]{d9ead3}Y lambda expr. &\cellcolor[HTML]{b7e1cd}Y LINQ + mapping &\cellcolor[HTML]{b7e1cd}Y LINQ + mapping &\cellcolor[HTML]{b7e1cd}Y LINQ + mapping &\cellcolor[HTML]{b7e1cd}Y LINQ + mapping \\
E1\_ColumnSorting &Y raw SQL &Y raw SQL &\cellcolor[HTML]{d9ead3}Y lambda expr. &\cellcolor[HTML]{b7e1cd}Y LINQ &\cellcolor[HTML]{b7e1cd}Y LINQ &\cellcolor[HTML]{b7e1cd}Y LINQ &\cellcolor[HTML]{b7e1cd}Y LINQ \\
E2\_Distinct &Y raw SQL &Y raw SQL &Y raw SQL &\cellcolor[HTML]{b7e1cd}Y LINQ &\cellcolor[HTML]{b7e1cd}Y LINQ &\cellcolor[HTML]{b7e1cd}Y LINQ &\cellcolor[HTML]{b7e1cd}Y LINQ \\
F1\_NestedJSONQuery &\cellcolor[HTML]{f4cccc}N custom JSON parsing &\cellcolor[HTML]{f4cccc}N custom JSON parsing &\cellcolor[HTML]{f4cccc}N custom JSON parsing &\cellcolor[HTML]{f4cccc}N custom converter &\cellcolor[HTML]{f4cccc}N custom JSON parsing &\cellcolor[HTML]{f4cccc}N custom JSON parsing &\cellcolor[HTML]{b7e1cd}Y LINQ \\
F2\_JSONArrayQuery &\cellcolor[HTML]{f4cccc}N custom JSON parsing &\cellcolor[HTML]{f4cccc}N custom JSON parsing &\cellcolor[HTML]{f4cccc}N custom JSON parsing &\cellcolor[HTML]{f4cccc}N custom converter &\cellcolor[HTML]{f4cccc}N custom JSON parsing &\cellcolor[HTML]{f4cccc}N custom JSON parsing &\cellcolor[HTML]{b7e1cd}Y LINQ \\
G1\_Union &Y raw SQL &Y raw SQL &Y raw SQL &\cellcolor[HTML]{b7e1cd}Y LINQ &\cellcolor[HTML]{b7e1cd}Y LINQ &\cellcolor[HTML]{b7e1cd}Y LINQ &\cellcolor[HTML]{b7e1cd}Y LINQ \\
G2\_Intersection &Y raw SQL &Y raw SQL &Y raw SQL &\cellcolor[HTML]{b7e1cd}Y LINQ &\cellcolor[HTML]{b7e1cd}Y LINQ &\cellcolor[HTML]{b7e1cd}Y LINQ &\cellcolor[HTML]{b7e1cd}Y LINQ \\
H1\_Metadata &Y raw SQL &Y raw SQL &Y raw SQL &Y raw SQL &Y raw SQL &Y raw SQL &Y raw SQL \\
\bottomrule
\end{tabular}
\end{threeparttable}
\end{table}
%\end{adjustwidth}
\end{landscape}
}



% \begin{center}
%     \captionof{table}{A1\_EntityIdenticalToTable}
%     \sisetup{table-format=4.1, separate-uncertainty}
%     \begin{tabular}{l S S S S S}
%         \hline
%         ORM & {Mean (\si{\micro\second})} & {Allocated (KB)} & {Gen0} & {Gen1} & {Gen2} \\
%         \hline
%         Dapper    & 750.2  & 6.24  & 0.0    & 0.0    & 0.0    \\
%         PetaPoco  & 721.8  & 7.99  & 0.9766 & 0.0    & 0.0    \\
%         RepoDB    & 747.5  & 8.77  & 0.9766 & 0.9766 & 0.0    \\
%         Linq2db   & 839.2  & 12.74 & 0.9766 & 0.0    & 0.0    \\
%         NHibernate & 742.1 & 36.06 & 3.9063 & 0.0    & 0.0    \\
%         EF6       & 854.1  & 90.43 & 9.7656 & 0.0    & 0.0    \\
%         EFCore    & 809.1  & 14.49 & 0.9766 & 0.0    & 0.0    \\
%         \hline
%     \end{tabular}
% \end{center}


\clearpage
\begin{landscape}
\begin{table}
\centering
\caption{Performance benchmark results - Mean time (\unit{\micro\second})}
\label{tab:benchmark_results_time}
\scriptsize
\def\arraystretch{1.35}
\begin{tabular}{
>{\raggedright\arraybackslash}p{50.00mm}
>{\raggedleft\arraybackslash}p{20.00mm}
>{\raggedleft\arraybackslash}p{20.00mm}
>{\raggedleft\arraybackslash}p{20.00mm}
>{\raggedleft\arraybackslash}p{20.00mm}
>{\raggedleft\arraybackslash}p{20.00mm}
>{\raggedleft\arraybackslash}p{20.00mm}
>{\raggedleft\arraybackslash}p{20.00mm}
>{\raggedleft\arraybackslash}p{20.00mm}
}
\toprule
\textbf{Namespace} &    \textbf{Dapper} &  \textbf{PetaPoco} &    \textbf{RepoDB} &   \textbf{Linq2db} &  \textbf{NHibernate} &        \textbf{EF6} &     \textbf{EFCore} \\
\midrule
A1\_EntityIdenticalToTable        &      750 &      722 &      748 &      839 &        742 &        854 &        809 \\
A2\_LimitedEntity                 &      718 &      726 &      735 &      858 &        958 &        882 &        799 \\
A3\_MultipleEntitiesFromOneResult &      736 &      736 &      745 &      898 &      1 001 &        945 &        832 \\
A4\_StoredProcedureToEntity       &  523 239 &  505 967 &  523 242 &  519 491 &    621 739 &    514 077 &    521 573 \\
B1\_SelectionOverIndexedColumn    &      744 &      759 &      762 &      868 &        817 &        905 &        806 \\
B2\_SelectionOverNonIndexedColumn &   43 612 &   41 100 &   43 052 &   43 708 &     86 419 &    125 169 &     48 016 \\
B3\_RangeQuery                    &   30 393 &   29 935 &   30 746 &   21 909 &     22 289 &     22 823 &     21 102 \\
B4\_InQuery                       &      856 &      855 &      876 &      954 &        945 &      1 408 &      1 147 \\
B5\_TextSearch                    &  747 728 &  747 721 &  745 551 &  745 319 &    746 464 &    746 264 &    744 754 \\
B6\_PagingQuery                   &    1 327 &    1 314 &    1 556 &    1 458 &      1 447 &      1 710 &      1 406 \\
C1\_AggregationCount              &   35 006 &   35 289 &  377 349 &   35 372 &     35 553 &     35 364 &     35 286 \\
C2\_AggregationMax                &    1 240 &    1 227 &    1 264 &    1 380 &      1 291 &      1 366 &      1 304 \\
C3\_AggregationSum                &   86 499 &   85 944 &   86 899 &   86 145 &     75 137 &     69 827 &     87 825 \\
D1\_OneToManyRelationship         &      781 &      770 &      806 &    2 997 &      1 042 &      1 581 &        895 \\
D2\_ManyToManyRelationship        &    4 584 &    4 480 &    6 031 &    9 383 &      5 940 &      6 913 &     14 964 \\
D3\_OptionalRelationship          &  187 693 &  172 030 &  393 631 &   91 326 &    533 489 &  1 993 146 &  2 548 411 \\
E1\_ColumnSorting                 &    4 725 &    4 622 &    4 963 &    4 845 &      6 316 &      6 204 &      4 885 \\
E2\_Distinct                      &    2 181 &    2 148 &    2 180 &    2 374 &      2 229 &      2 340 &      2 270 \\
F1\_JSONObjectQuery               &    1 459 &    1 478 &    1 457 &    1 552 &      1 494 &      1 490 &      1 544 \\
F2\_JSONArrayQuery                &    1 723 &    1 718 &    1 732 &    1 795 &      1 761 &      1 771 &      1 839 \\
G1\_Union                         &      728 &      715 &      723 &    1 580 &      1 538 &      1 613 &      1 561 \\
G2\_Intersection                  &      739 &      718 &      730 &    1 596 &      1 534 &      1 622 &      1 566 \\
H1\_Metadata                      &      851 &      843 &      844 &      925 &        866 &        899 &        814 \\
\bottomrule
\end{tabular}
\end{table}
\end{landscape}

\clearpage
\begin{landscape}
\begin{table}
\centering
\caption{Performance benchmark results - Allocated memory (KB)}
\label{tab:benchmark_results_memory}
\scriptsize
\def\arraystretch{1.35}
\begin{tabular}{
>{\raggedright\arraybackslash}p{50.00mm}
>{\raggedleft\arraybackslash}p{20.00mm}
>{\raggedleft\arraybackslash}p{20.00mm}
>{\raggedleft\arraybackslash}p{20.00mm}
>{\raggedleft\arraybackslash}p{20.00mm}
>{\raggedleft\arraybackslash}p{20.00mm}
>{\raggedleft\arraybackslash}p{20.00mm}
>{\raggedleft\arraybackslash}p{20.00mm}
>{\raggedleft\arraybackslash}p{20.00mm}
}
\toprule
\textbf{Namespace}                         &    \textbf{Dapper} &  \textbf{PetaPoco} &    \textbf{RepoDB} &   \textbf{Linq2db} & \textbf{NHibernate} &      \textbf{EF6} &    \textbf{EFCore} \\
\midrule
A1\_EntityIdenticalToTable        &      6.24 &     7.99 &      8.77 &     12.74 &      36.06 &    90.43 &    14.49 \\
A2\_LimitedEntity                 &      4.74 &     9.64 &      7.99 &     18.25 &      50.70 &   115.18 &    16.11 \\
A3\_MultipleEntitiesFromOneResult &      7.28 &    13.03 &     10.38 &     27.74 &      84.00 &   211.61 &    25.85 \\
A4\_StoredProcedureToEntity       &  24294.84 & 16490.15 &  17007.14 &  16490.93 &   91569.96 & 35270.68 & 29008.69 \\
B1\_SelectionOverIndexedColumn    &      8.63 &    14.30 &      9.79 &     14.63 &      49.46 &   102.35 &    12.88 \\
B2\_SelectionOverNonIndexedColumn &   7112.75 &  3308.54 &   3410.07 &   3314.16 &   21371.02 & 50014.45 &  5847.86 \\
B3\_RangeQuery                    &    989.93 &   468.38 &    480.34 &    475.55 &    3048.82 &  6747.98 &   819.92 \\
B4\_InQuery                       &     16.92 &    20.34 &     16.64 &     19.58 &      87.42 &   354.18 &    19.85 \\
B5\_TextSearch                    &   1169.29 &   587.24 &    599.41 &    592.78 &    3472.99 &  8113.77 &   979.42 \\
B6\_PagingQuery                   &     32.59 &    25.06 &     18.25 &     26.58 &     128.29 &   265.26 &    40.52 \\
C1\_AggregationCount              &      2.75 &     8.79 &  59994.30 &     12.46 &      41.75 &   117.99 &    12.04 \\
C2\_AggregationMax                &      1.97 &     4.19 &      1.77 &      8.12 &      23.06 &    77.09 &     5.69 \\
C3\_AggregationSum                &      2.09 &     4.41 &      1.35 &      9.25 &      25.31 &    81.72 &     6.85 \\
D1\_OneToManyRelationship         &     16.46 &    17.94 &     20.23 &     24.58 &      88.71 &   115.27 &    26.44 \\
D2\_ManyToManyRelationship        &    433.35 &   381.13 &    854.28 &    352.15 &    1162.43 &  1395.95 &  2555.79 \\
D3\_OptionalRelationship          &  40923.17 & 24909.64 &  20945.39 &   9083.71 &  144094.98 &278566.55 &175866.23 \\
E1\_ColumnSorting                 &    349.40 &   156.96 &    162.62 &    162.54 &    1354.02 &  2000.95 &   348.07 \\
E2\_Distinct                      &      2.42 &     4.78 &      1.66 &      9.34 &      23.72 &    76.19 &     8.54 \\
F1\_JSONObjectQuery               &     21.71 &    45.25 &     31.82 &     27.09 &      86.13 &    73.38 &    51.68 \\
F2\_JSONArrayQuery                &      6.04 &    18.12 &      8.25 &     11.30 &      74.31 &    57.24 &    16.86 \\
G1\_Union                         &      2.39 &     3.45 &      1.40 &     18.03 &      60.07 &   136.10 &    20.35 \\
G2\_Intersection                  &      2.17 &     3.25 &      1.30 &     17.91 &      61.36 &   135.98 &    21.73 \\
H1\_Metadata                      &      1.98 &     5.10 &      1.21 &      5.61 &      35.53 &    46.52 &    10.99 \\
\bottomrule
\end{tabular}
\end{table}
\end{landscape}

\chapter{Temporary section for notes}

\section{Abstraction}
\begin{verbatim}
- entity
	- name
	- properties[]
		- name
		- data type
		- is nullable?
		- is primary key?
		- constraints[] (maxlength, range, ...)
	- relationships[]
		- multiplicity
		- target entity
		- navigation property
	- inheritance?
- mapping
	- entity
	- mapped properties[]
		- property name
		- column name
		- database data type
		- other details[]
	- relationship mapping[]
		- type
		- target entity
		- navigation property
		- join table?
- query
	- source (context,session)
	- result variable
		- data type
		- name
	- select[] (can have complex transformations in linq)
		- columns
		- type
		- name
	- where[]
		- expression
	- joins[]
		- join type
		- source entity
		- target entity
		- on condition
	- order by[]
	- group by[]
	- having[]
	- pagination
		- skip
		- take
	- parameters
		- name
		- type
		- value
	- execution
		- operator (ToList, First, Single)
		- default if empty?
\end{verbatim}


\chapter*{Conclusion}
\addcontentsline{toc}{chapter}{Conclusion}


%%% Bibliography
%%% Bibliography (literature used as a source)
%%%
%%% We employ biblatex to construct the bibliography. It processes
%%% citations in the text (e.g., the \cite{...} macro) and looks up
%%% relevant entries in the bibliography.bib file.
%%%
%%% See also biblatex settings in thesis.tex.

%%% Generate the bibliography. Beware that if you cited no works,
%%% the empty list will be omitted completely.

% We let bibliography items stick out of the right margin a little
\def\bibfont{\hfuzz=2pt}

\printbibliography[heading=bibintoc]

%%% If case you prefer to write the bibliography manually (without biblatex),
%%% you can use the following. Please follow the ISO 690 standard and
%%% citation conventions of your field of research.

% \begin{thebibliography}{99}
%
% \bibitem{lamport94}
%   {\sc Lamport,} Leslie.
%   \emph{\LaTeX: A Document Preparation System}.
%   2nd edition.
%   Massachusetts: Addison Wesley, 1994.
%   ISBN 0-201-52983-1.
%
% \end{thebibliography}


%%% Figures used in the thesis (consider if this is needed)
\listoffigures

%%% Tables used in the thesis (consider if this is needed)
%%% In mathematical theses, it could be better to move the list of tables to the beginning of the thesis.
\listoftables

%%% Abbreviations used in the thesis, if any, including their explanation
%%% In mathematical theses, it could be better to move the list of abbreviations to the beginning of the thesis.
\chapwithtoc{List of Abbreviations}

\begin{table}[h!]
    \centering

    \begin{tabular}{ll}
        \toprule
        abbreviation & meaning                        \\
        \midrule
        ORM     & Object-relational mapping           \\
        EF      & Entity Framework                    \\
        ID      & identifier                          \\
        MSSQL   & Microsoft SQL Server                \\
        LINQ    & Language Integrated Query           \\
        WWI     & World Wide Importers dataset (\autoref{sec:database})       \\
                &                                     \\
        \bottomrule
    \end{tabular}

    \caption{List of abbreviations}

\end{table}

%%% Doctoral theses must contain a list of author's publications
\ifx\ThesisType\TypePhD
\chapwithtoc{List of Publications}
\fi

%%% Attachments to the thesis, if any. Each attachment must be referred to
%%% at least once from the text of the thesis. Attachments are numbered.
%%%
%%% The printed version should preferably contain attachments, which can be
%%% read (additional tables and charts, supplementary text, examples of
%%% program output, etc.). The electronic version is more suited for attachments
%%% which will likely be used in an electronic form rather than read (program
%%% source code, data files, interactive charts, etc.). Electronic attachments
%%% should be uploaded to SIS. Allowed file formats are specified in provision
%%% of the rector no. 72/2017. Exceptions can be approved by faculty's coordinator.
\appendix
\chapter{Attachments}

\section{First Attachment}

\end{document}
